\section{Threats to Validity}\label{sec:threats}

This section lists the threats known before any confirmatory data. Threats that the results reveal
belong to Section~\ref{sec:discussion}.

\subsection{Scope}\label{sec:threat-scope}

Every claim is about one server, as built, on two hosts, over loopback. A network between client
and server would add latency and interrupt work that loopback does not have, and both would dilute
the cost of detection further. The cost claims hold only for load of one protocol at a time; the
mixed-protocol cell is secondary. Each claim is made per protocol, backend and host, never pooled.

The TLS cells of the cost family have little power to show the cost of detection, since the
handshake dominates each connection; they are not read as evidence about it. In the SSH cells the
server speaks at accept in dedicated mode, but only after the client's first bytes in one-port
mode. So the two arms do not exchange messages in the same order. M1 on IOCP compares replay with
a peek that switches to replay at its first undecided read, because Windows offers no low-water
mark. It measures the limits of the interface as much as the idea.

\subsection{Engineering before the Freeze}\label{sec:threat-engineering}

The server was engineered until it won where it could, against the proxies and libraries, on
development data, before the code freeze. Those runs could not time one-port mode against
dedicated mode. The A/A pilot that sizes the cost family holds no one-port data. So neither the
margin nor \RC{} could be tuned to a cost contrast. The design choices of the relay and of the
footprint test were informed by development runs against the competitors. The confirmatory runs
are new data on the frozen binary, but a reader should weigh that history.

\subsection{The Footprint Test}\label{sec:threat-footprint}

$W$ is the footprint of each system as configured, not of its detection design alone. Two
residuals of the socket-buffer accounting remain, and both sit with a system that leaves bytes
queued. Such systems are nginx and Envoy in the partial-ClientHello case, through their peek
paths, and the server only where its detection mode is peek. The first raises that system's $W$ and the second lowers
it. In the silent case neither applies.

On L one socket-buffer cache is merged with two caches that hold neither sockets nor socket
buffers. Their growth is subtracted with the cache's. Growth caused by the host favours neither
arm, since each session runs $XYYX$ with $X$ drawn, and a drift that is linear within a session
falls on both arms alike. Growth that a system causes itself goes with its arm, and randomisation
does not remove it. Three things bound it, though none is a test. Each window's processes start
and pass the probe before the baseline. The co-tenants grow only when the kernel creates workqueue
pools or chained scatter-gather lists. And the shared cache's growth is reported per arm beside
every cell. In the silent case no system leaves a byte queued, so that growth is read directly.

The server's $W$ is at least its own $K_s$, most of which is the memory of the two loopback socket
ends that \Kbase{} measures. Against a competitor whose own footprint is close to a bare socket's,
no change to the server can reach the margin of B3. Such a cell is reported as not shown.

The JVM and Go systems hold memory that their collectors manage. Each is collected before every
reading, and its cells are descriptive.

\subsection{Windows}\label{sec:threat-w}

W's frozen procedure names a CPU frequency counter, tested against a known load before the code
freeze. If no counter passes that test, W's windows are validated without a frequency rule: a
change of clock speed within a W session would then go undetected.
\ForAlex{The first test, in M6b, did not pass: under its control the counter did not move, and
neither did a spinner's work rate (design/status-m6b.md). A retest with a capped plan is prepared.
State the outcome here once it is in the revision log.}

On W the server's CPU time comes from the process times that Windows reports. Windows may charge
that time in whole clock ticks, of \WTickMs{}~ms on W, to the thread that runs at each tick. If it
does, then at a low server load the change over a window is a sample, not a sum. This would affect
only WL4 on W, a secondary metric.
\ForAlex{M6c's development runs on W saw exactly this: every server CPU time of the W A/A job was a
whole multiple of the tick (design/status-m6b.md on branch \texttt{m6b-windows}). The W runner
there (commit 485d7ae, not merged into main) also records the server's processor cycles. Which
value is W's WL4 is open for the coordinator.}

\ForAlex{Two more readings of M6b are not yet in the revision log: on W, opgen sets
\texttt{SO\_REUSE\_UNICASTPORT} in place of \texttt{IP\_BIND\_ADDRESS\_NO\_PORT}, and W has no
counter of listen overflows, so a refused connection shows only as a failed connect. M6c also found
that W's churn of the protocols whose client closes first may be bound by the generator. Decide
which of these the paper states.}

\subsection{Cells at Risk of Not Running}\label{sec:threat-cells}

In M2's TLS cells, the back end that terminates TLS may be the relay arm's limit in closed loop.
Its cores would then run busier than the rule for back-end cores allows, every rate session would
be invalid, and the cell would have no rate. By the frozen rule such a cell runs no window and
enters Holm with \PIsOne{}, and M2's claim narrows to HTTP/1.1.
\ForAlex{The coordinator's decision on M2's TLS cells is open (design/status.md, M7d).}

TLS with session resumption is not run, for the reason given in Section~\ref{sec:secondary}.

sslh-ev version \SslhVersion{} sets up its event watcher for a descriptor number only the first
time it sees that number~\cite{sslh2026}. A connection accepted onto a number closed in the same
pass of its loop can then go unread. Its churn windows may fail the error rule, and its cells of M3
would enter Holm with \PIsOne{}. The narrowing order then applies.

\subsection{Competitors and Instruments}\label{sec:threat-instruments}

Each competitor runs the configuration its project documents for the compared feature. A better
configuration may exist that we did not find. HAProxy's client, server and connect timeouts stay
at its default, which is none. Envoy's M3 and B3 configurations hold the TLS inspector and not the
HTTP inspector, since no route of theirs reads HTTP.

The pre-specification cites Linux source lines at tag \KernelTag. The runs use kernel \KernelL. The
revision log records that each cited line is unchanged in content there; the distribution's own
patches on top of the tag were not compared.

The generator's closed loop holds \ConnSlots{} connections per server core, so a closed-loop cell
measures throughput at that concurrency only. The open loop measures the time to first byte from
each exchange's due time, so a stall of the generator shows as latency rather than hiding.
